\documentclass[11pt]{article}
\pdfinfoomitdate=1
\pdftrailerid{}
\pdfsuppressptexinfo=15
\usepackage[margin=1in]{geometry}
\usepackage{amsmath,amssymb,amsthm,mathtools,booktabs,longtable,array,microtype}
\usepackage[T1]{fontenc}
\usepackage{lmodern}
\usepackage{xcolor}
\usepackage[colorlinks=true,linkcolor=blue!55!black,citecolor=blue!55!black,urlcolor=blue!55!black]{hyperref}
\newtheorem{theorem}{Theorem}[section]
\newtheorem{lemma}[theorem]{Lemma}
\newtheorem{proposition}[theorem]{Proposition}
\newtheorem{corollary}[theorem]{Corollary}
\newtheorem{definition}[theorem]{Definition}
\newtheorem{remark}[theorem]{Remark}
\newtheorem{firewall}[theorem]{Firewall}
\newcommand{\R}{\mathbb R}
\newcommand{\N}{\mathbb N}
\newcommand{\Acal}{\mathcal A}
\newcommand{\Rcal}{\mathcal R}
\newcommand{\Fcal}{\mathcal F}
\newcommand{\Scal}{\mathcal S}
\newcommand{\Lcal}{\mathcal L}
\newcommand{\CSHT}{\mathrm{CSHT}_{67}}
\newcommand{\LAPBR}{\mathrm{LAPBR}_{67}}
\title{Threshold-Flexible Causality for the Factor--67 Annular Scalar:\\Exact Source Completion, a Subcritical-Depth No-Go, and the Critical-Saddle Frontier}
\author{Agentic Polymath Project}
\date{August 18, 2026}
\begin{document}
\maketitle

\begin{abstract}
We reconstruct the threshold-flexible factor--67 causal method at the level of
individual native rough occurrences.  Every rough edge carries its literal
coefficient $r=p^{-1/2}$, and every threshold split is the one-use partition
$(r-a)+a=r$.  The missing one-prime compensation $r-2r^2$, and in general the
nonlocal resolvent term $(R-A)F$, are therefore compulsory.  We prove that the
root observation of the completed owner tree is the native annular scalar, not
the finite $P_{61}$ base; at $X=184$ the distinction exceeds $1.36314$.

The source ledger does not imply positivity.  Let $C_L(X)$ be the exact native
current obtained by stopping before rough depth $L$, with $L$ even.  We prove
that
\[
 L(X)\log L(X)=o(\log\log X)
 \quad\Longrightarrow\quad C_L(X)<0
\]
for all sufficiently large $X$.  This refutes the large-prime residual theorem
$\LAPBR$ proposed for the adaptive depth $L=O(\log\log\log X)$.  A directed
256-bit calculation also certifies the adjacent finite sign change
$C_2(32605)>0>C_2(32606)$ and $C_2(61841)<-21.3$, while the native scalar at
$61841$ is positive.  On the constructive side, we prove positivity whenever
the least rough prime is at least $Y^\theta$ with $\theta>e^{-1}$.

The surviving producer must operate in the critical rough-depth window
$k\asymp\log\log X$ and transport complete odd-depth native occurrences into
even-depth occurrences with one common source ledger.  We formulate this as
$\CSHT$.  The paper proves the source and no-go theorems unconditionally; it
does not prove $\CSHT$ or the Riemann Hypothesis.
\end{abstract}

\begin{center}
\fbox{\parbox{0.92\textwidth}{\textbf{Scientific status.}
The threshold-flexible subcritical-depth method is decided negatively.  The
new exact source theorems and high-prime sector are unconditional.  The
critical-saddle transport and the Riemann Hypothesis remain unproved.}}
\end{center}

\tableofcontents

\section{The annular scalar and the frozen finite base}

For $j\ge2$, let $c_X(j)$ denote the native Mobius component row used by the
fixed-row Mellin consumer.  We study the annular $5{:}3$ scalar
\begin{equation}\label{eq:native-scalar}
 \Acal_X
 =5\bigl(c_X(2)-c_{X/4}(2)\bigr)
 +3\bigl(c_X(3)-c_{X/4}(3)\bigr).
\end{equation}
Before Mobius convolution its coefficient dictionary is
\begin{equation}\label{eq:qstar}
 q_*(2)=15,\quad q_*(3)=6,\quad q_*(4)=3,
 \quad q_*(n)=6\quad(n\ge5),
\end{equation}
with $q_*(1)=0$.

Put $P_{61}=\prod_{p\le61}p$.  The grouped finite-Euler annular scalar is
\begin{equation}\label{eq:bdef}
 b(X)=F_{61}(X)
 =\sum_{d\mid P_{61}}{\mu(d)\over\sqrt d}A_*(X/d),
\end{equation}
where
\[
 A_*(Y)=\sum_{n\ge2}{q_*(n)\over\sqrt n}
        \min\!\left(\log4,\log{Y\over n}\right)_+.
\]
We use the following independently certified theorem as a frozen arithmetic
input.

\begin{proposition}[Sharp finite-$P_{61}$ base]\label{prop:base}
The function $b$ is nonnegative for every real $Y\ge1$ and
\[
 b(Y)=a\sqrt Y+O(1),
 \qquad
 a=12\prod_{p\le61}\left(1-{1\over p}\right)>0.
\]
For $Y\ge67$, its unsigned companion $M(Y)$ satisfies
\[
 {1\over42}M(Y)\le b(Y)\le {1\over8}M(Y).
\]
The global ratio minimum occurs at $Y=184$ and is
\[
 {b(184)\over M(184)}
 =0.02398429355876630467327315865\ldots>{1\over42}.
\]
\end{proposition}

The directed proof of Proposition~\ref{prop:base} is part of the frozen base
packet.  The present paper does not reuse the false earlier constant $1/40$.

\section{The native root and one exact owner per rough occurrence}

Let
\[
 \Rcal_{67}^{\mathrm{sf}}
 =\{m:m\text{ is squarefree and }p\mid m\Rightarrow p\ge67\}.
\]

\begin{theorem}[Native rough convolution]\label{thm:native-root}
For every real endpoint $X$,
\begin{equation}\label{eq:root-convolution}
 \boxed{
 \Acal_X=\sum_{m\in\Rcal_{67}^{\mathrm{sf}}}
 {\mu(m)\over\sqrt m}\,b(X/m).
 }
\end{equation}
The sum is finite.  If $m=p_1\cdots p_t$ with
$67\le p_1<\cdots<p_t$, then the unique history owner of that occurrence has
coefficient $m^{-1/2}$ and parity $(-1)^t$.
\end{theorem}

\begin{proof}
Every squarefree native index has a unique factorization
\[
 k=d p_1\cdots p_t,
 \qquad d\mid P_{61},
 \qquad67\le p_1<\cdots<p_t.
\]
The finite and rough Euler functions have disjoint prime support, hence
\[
 \mu=\mu_{P_{61}}*\mu_{\Rcal_{67}^{\mathrm{sf}}}.
\]
Convolving \eqref{eq:qstar} and applying the same annular hinge proves
\eqref{eq:root-convolution} coefficientwise.

Moving one prime $p$ into the ordered history replaces $(X,k)$ by
$(X/p,k/p)$.  Both native invariants are exact:
\[
 {X/p\over k/p}={X\over k},
 \qquad
 p^{-1/2}(k/p)^{-1/2}=k^{-1/2}.
\]
The paired parity channel swaps once, so the complete sign is
$\mu(d)(-1)^t=\mu(k)$.  Unique factorization gives one ordered first-owner
history.  The hinge makes the sum finite.
\end{proof}

\begin{firewall}[Ownership is not positivity]
Theorem~\ref{thm:native-root} proves coefficient, activation and parity
provenance.  It does not imply that a stopped current or a terminal leaf has
nonnegative signed observation.
\end{firewall}

At $X=184$, only the rough primes $67,71,73,79,83,89$ occur.  Directly from
\eqref{eq:root-convolution},
\[
 \Acal_{184}=b(184)-D_{184},
 \qquad
 D_{184}>1.3631478826704517.
\]
The directed replay gives
\[
 b(184)\in[10.69357964877382966,10.69357964877383023],
\]
\[
 \Acal_{184}\in[9.33043176610337787,9.33043176610337847].
\]
Thus any low-child recombination returning $b(184)$ fails the native-root
identity before positivity is discussed.

\section{Threshold policies and the exact residual}

Let the finite rough-history DAG be ordered from parents to children.  Let $R$
be its raw nonnegative edge operator; the entry on a prime-$p$ edge is
$r=p^{-1/2}$.  Let $S$ denote parity-channel swap.

\begin{theorem}[One-use threshold ledger]\label{thm:threshold-ledger}
At an edge $v\to w$, choose any $a_{v,w}$ with
$0\le a_{v,w}\le r_{v,w}$.  The raw positive paired-source occurrence has the
unique one-use partition
\begin{equation}\label{eq:edge-split}
 r_{v,w}SP_w
 =(r_{v,w}-a_{v,w})SP_w\oplus a_{v,w}SP_w.
\end{equation}
Consequently every native occurrence is either assigned one deterministic
owner or is partitioned into disjoint suboccurrences whose coefficient sum is
its original coefficient.  No threshold policy changes the raw coefficient.
\end{theorem}

\begin{proof}
Equation \eqref{eq:edge-split} is an identity in the positive paired-source
cone.  Induct along the unique ordered history of
Theorem~\ref{thm:native-root}.  At every edge the two subcoefficients add to
$r$.  Multiplication along the path therefore gives exactly
$\prod_i p_i^{-1/2}$.  Distinct owner labels are disjoint by construction.
\end{proof}

Let $b$ now denote the vector of local grouped observations and let $F$ denote
the exact native completed-parity scalar.  The raw recursion is
\begin{equation}\label{eq:raw-recursion}
 F+RF=b.
\end{equation}

\begin{theorem}[Raw-to-contracted resolvent]\label{thm:resolvent}
For any nonnegative threshold operator $A$ with $0\le A\le R$, the same native
scalar satisfies
\begin{equation}\label{eq:contracted-recursion}
 F+AF=g
\end{equation}
if and only if
\begin{equation}\label{eq:forced-current}
 \boxed{g=(I+A)(I+R)^{-1}b=b-(R-A)F.}
\end{equation}
The positive current source is
\begin{equation}\label{eq:positive-current}
 C_v=B_v\oplus\bigoplus_w(r_{v,w}-a_{v,w})SP_w,
\end{equation}
while the recursive source is
$\bigoplus_w a_{v,w}SP_w$.
\end{theorem}

\begin{proof}
Since $R$ is nilpotent, $F=(I+R)^{-1}b$.  Substitution into
\eqref{eq:contracted-recursion} gives \eqref{eq:forced-current}.  Equation
\eqref{eq:positive-current} is the direct sum of the residual pieces in
\eqref{eq:edge-split}.  Adding the recursive pieces recovers the raw parent.
Signed observation gives exactly \eqref{eq:forced-current}.
\end{proof}

\subsection{The one-prime correction}

For one prime, let $r=p^{-1/2}$.  If a proposed packet has already charged one
$r^2$ child in the current and another $r^2$ child recursively, the omitted
positive source is
\begin{equation}\label{eq:delta}
 \boxed{\delta_p=r-2r^2=r(1-2r)>0.}
\end{equation}
Indeed $\delta_p+2r^2=r$.  A cleaner split with one recursive $r^2$ child has
residual $r-r^2$.  Both are instances of $(R-A)F$.

\begin{corollary}[Raw-exposure conservation]\label{cor:exposure}
For every child scalar $F_w$,
\[
 (r-a)F_w+aF_w=rF_w.
\]
Threshold contraction relocates raw exposure; it does not reduce it.
\end{corollary}

\begin{firewall}[Minimal two-node separator]\label{fw:two-node}
Take $b=(0,1)$, one raw edge $r>0$, and one threshold $0\le a<r$.  Then
$F=(-r,1)$ and the exact parent current is
\[
 g_0=F_0+aF_1=a-r<0.
\]
Thus no universal local one-channel positive current can implement a strict
threshold contraction.
\end{firewall}

\section{The odd-history stress case}

The preceding source identities preserve parity rather than hiding it.  This is
mandatory at
\[
 X=67\cdot71\cdot13=61841,
 \qquad h=(67),
 \qquad(p,y)=(71,13).
\]
A 192-bit exact dyadic enclosure over all 239 active $P_{61}$ divisors proves
\begin{equation}\label{eq:odd-gap}
 E_T(71,13)-O_T(71,13)>17.
\end{equation}
The incoming history is odd, so a leafwise canonical Hall operation would need
the reverse inequality $O_T\ge E_T$.  Equation \eqref{eq:odd-gap} excludes it.
The threshold ledger passes this test because it never reorients the child; it
also shows why ownership alone is insufficient.

\section{Exact stopped currents}

For even $L\ge2$, define the natural current through depth $L-1$ by
\begin{equation}\label{eq:CL}
 C_L(X)=
 \sum_{\substack{m\in\Rcal_{67}^{\mathrm{sf}}\\\omega(m)<L}}
 {\mu(m)\over\sqrt m}\,b(X/m).
\end{equation}
This is the literal current obtained by expanding the raw recursion for
$L-1$ levels.  Since $L$ is even, the unresolved depth-$L$ frontier has canonical
parity, but that fact does not determine the sign of $C_L$.

The finite replay gives
\[
 C_2(32605)>0>C_2(32606),
\]
with strict 256-bit intervals, and
\[
 C_2(61841)<-21.3,
 \qquad
 \Acal_{61841}>9.5.
\]
Thus even an exact negative current can telescope with its remaining frontier
to a positive native root.

\section{Subcritical threshold depths are eventually negative}

\begin{theorem}[Subcritical-depth no-go]\label{thm:no-go}
Let $L=L(X)\ge2$ be even and suppose
\[
 L(X)\log L(X)=o(\log\log X).
\]
Then the exact current \eqref{eq:CL} satisfies
\[
 \boxed{C_L(X)<0}
\]
for every sufficiently large $X$.
\end{theorem}

\begin{proof}
Put $k=L-1$, so $k$ is odd, and define
\[
 z_X=\sum_{67\le p\le X}{1\over p},
 \qquad
 W=X^{1/(2k)},
 \qquad
 z_W=\sum_{67\le p\le W}{1\over p}.
\]
By Proposition~\ref{prop:base}, there are constants $0<c<C$ such that
\[
 b(Y)\ge c\sqrt Y\quad(Y\text{ sufficiently large}),
 \qquad
 |b(Y)|\le C\sqrt Y\quad(Y\ge1).
\]
Every product of $k$ primes at most $W$ is at most $\sqrt X$, so all selected
last-layer arguments are at least $\sqrt X$.  If $e_j(S)$ denotes the elementary
symmetric sum of degree $j$ in $\{1/p:p\in S\}$, then we may discard all other
negative levels and obtain
\begin{equation}\label{eq:upper-current}
 C_L(X)
 \le C\sqrt X\sum_{j=0}^{k-1}e_j(\{67\le p\le X\})
 -c\sqrt X\,e_k(\{67\le p\le W\}).
\end{equation}

The multinomial expansion gives
\[
 e_j(\{p\le X\})\le {z_X^j\over j!}.
\]
Let $s_2=\sum_{p\ge67}p^{-2}<\infty$.  In the ordered expansion of $z_W^k$,
the total weight of tuples with a repeated coordinate is at most
$\binom{k}{2}s_2z_W^{k-2}$.  Hence
\begin{equation}\label{eq:eklower}
 e_k(\{p\le W\})
 \ge {z_W^k\over k!}
 \left(1-{\binom{k}{2}s_2\over z_W^2}\right).
\end{equation}

The hypothesis implies $k/z_X\to0$, so
\[
 \sum_{j=0}^{k-1}{z_X^j\over j!}
 \le(1+o(1)){z_X^{k-1}\over(k-1)!}.
\]
Mertens' theorem gives
\[
 z_W=z_X-\log(2k)+o(1).
\]
Therefore
\begin{align*}
 {e_k(\{p\le W\})
  \over z_X^{k-1}/(k-1)!}
 &\ge(1-o(1)){z_W\over k}
 \left({z_W\over z_X}\right)^{k-1}\\
 &\longrightarrow\infty.
\end{align*}
Indeed $k\log k=o(z_X)$ makes the exponential factor tend to one, while
$z_W/k\to\infty$.  The negative last layer in
\eqref{eq:upper-current} dominates all preceding positive layers.  This proves
the theorem.
\end{proof}

\begin{remark}
Theorem~\ref{thm:no-go} includes every fixed even depth and every polylogarithmic
adaptive depth.  It does not claim that a stopping depth in the critical window
$L\asymp\log\log X$ fails.
\end{remark}

\section{The large-prime adaptive residual is false}

A recent adaptive proposal sets
\[
 Z_X=(\log X)^{1/4}
\]
and takes the smallest even $L_X$ satisfying
\[
 L_X\ge8\left(1+\sum_{67\le p\le Z_X}{1\over p}\right).
\]
It decomposes
\[
 C_{L_X}(X)=\Scal_X+\Lcal_X,
\]
where $\Scal_X$ uses only primes at most $Z_X$ and $\Lcal_X$ contains histories
meeting a larger prime.  The small-prime theorem proves $\Scal_X>0$ eventually.

\begin{corollary}[Refutation of $\LAPBR$]\label{cor:lapbr}
For the preceding schedule,
\[
 \boxed{\Lcal_X<0}
\]
for all sufficiently large $X$.
\end{corollary}

\begin{proof}
Mertens' theorem gives $L_X=O(\log\log\log X)$, hence
$L_X\log L_X=o(\log\log X)$.  Theorem~\ref{thm:no-go} gives
$C_{L_X}(X)<0$.  Since $\Scal_X>0$,
\[
 \Lcal_X=C_{L_X}(X)-\Scal_X<0.
\]
\end{proof}

Thus the proposed universal inequality $\LAPBR$ is not an open estimate waiting
for a sharper large sieve; it has the wrong sign asymptotically.

\section{A positive high-least-prime sector}

The negative result suggests separating genuinely terminal large primes from
the critical bulk.

\begin{theorem}[High-least-prime positivity]\label{thm:high-prime}
Fix $\theta>e^{-1}$.  Define
\[
 \Fcal_{p_0}(Y)=
 \sum_{\substack{m\in\Rcal_{67}^{\mathrm{sf}}\\P^-(m)\ge p_0}}
 {\mu(m)\over\sqrt m}b(Y/m).
\]
Uniformly for states satisfying $p_0\ge Y^\theta$,
\[
 \boxed{\Fcal_{p_0}(Y)>0}
\]
for all sufficiently large $Y$.
\end{theorem}

\begin{proof}
Since $3\theta>1$, an active rough product contains at most two primes.  Hence
\[
 \Fcal_{p_0}(Y)=b(Y)
 -\sum_{p\ge p_0}{1\over\sqrt p}b(Y/p)
 +\sum_{p<q\atop p,q\ge p_0}{1\over\sqrt{pq}}b(Y/(pq)).
\]
The last sum is nonnegative by Proposition~\ref{prop:base}.  Moreover,
\begin{align*}
 \sum_{Y^\theta\le p\le Y/2}{1\over\sqrt p}b(Y/p)
 &=a\sqrt Y\sum_{Y^\theta\le p\le Y/2}{1\over p}
   +O\left(\sum_{p\le Y}p^{-1/2}\right)\\
 &=a\sqrt Y\bigl(\log(1/\theta)+o(1)\bigr)
   +O(\sqrt Y/\log Y).
\end{align*}
Because $\log(1/\theta)<1$ and $b(Y)=a\sqrt Y+O(1)$, the one-prime subtraction
is strictly smaller than the parent for large $Y$.  Positivity follows.
\end{proof}

\section{Reinvented path: the critical-saddle transport}

Put
\[
 T_k(X)=
 \sum_{\substack{m\in\Rcal_{67}^{\mathrm{sf}}\\\omega(m)=k}}
 {1\over\sqrt m}b(X/m).
\]
Then
\begin{equation}\label{eq:depth-series}
 \Acal_X=\sum_{k\ge0}(-1)^kT_k(X).
\end{equation}
Theorem~\ref{thm:no-go} says that stopping substantially before the saddle of
the reciprocal-prime elementary symmetric sequence exposes a dominant odd
layer.  Theorem~\ref{thm:high-prime} removes a genuine high-least-prime sector.
The remaining mass lies near
\[
 k\asymp\sum_{p\le X}{1\over p}\asymp\log\log X.
\]

\begin{definition}[Critical-Saddle History Transport, $\CSHT$]
For every sufficiently large real $X$, remove the sector covered by
Theorem~\ref{thm:high-prime}.  On the remaining active history DAG, construct a
nonnegative one-use flow from every odd-depth native occurrence to even-depth
native occurrences such that:
\begin{enumerate}
\item every source coefficient, activation, first owner and history parity is
      preserved;
\item each even occurrence is used at most once;
\item all odd scalar demand is covered;
\item if target or physical coordinates are included, the same edge
      coefficient is used in every coordinate.
\end{enumerate}
Equivalently, produce the exact finite Hall/Lorenz primal or an exact separating
dual on every activation cell.
\end{definition}

\begin{theorem}[Conditional conclusion]\label{thm:conditional}
If $\CSHT$ holds, then $\Acal_X\ge0$ eventually.  The exact annular Mellin
transform is
\[
 (1-4^{-s})\left[
 {6\over s^2}
 -{3(1-2^{-z})(2-2^{-z})\over s^2\zeta(z)}
 \right],
 \qquad z=s+\tfrac12.
\]
The finite numerator is zero-free for $\Re z>0$.  The fixed-sign
Mellin--Landau argument would therefore imply RH.
\end{theorem}

\begin{proof}
The transport makes the even side of \eqref{eq:depth-series} dominate the odd
side.  The Mellin identity and zero-free finite numerator are exact.  Landau's
real-abscissa theorem gives the stated conditional implication.
\end{proof}

Theorem~\ref{thm:conditional} is not an unconditional RH proof because $\CSHT$
is open.  Unlike the refuted threshold scheme, it identifies the depth and the
source coordinates where a future proof must genuinely win.

\section{Replay and mandatory stress cases}

The deterministic packet contains an MPFR 4.2.2 producer at 256 bits and an
exact rational verifier.  It checks:
\begin{center}
\begin{tabular}{p{0.32\textwidth}p{0.58\textwidth}}
\toprule
Case & Certified contract\\
\midrule
one prime $p=67$ & $\delta_p=r-2r^2>0$ and both residual conventions sum to $r$\\
$X=184$ & finite $P_{61}$ base differs from the native root by more than $1.36$\\
$67\cdot71\cdot13$ & 239-atom odd-history target gap exceeds $17$\\
two-node thinning & exact current $a-r<0$\\
$X=32605,32606$ & directed adjacent depth-two sign change\\
$X=61841$ & depth-two current below $-21.3$ while native scalar exceeds $9.5$\\
raw exposure & $(r-a)+a=r$ edge by edge and through the resolvent\\
\bottomrule
\end{tabular}
\end{center}

Expected replay verdict:
\[
\texttt{PASS\_T97700\_THRESHOLD\_FLEXIBLE\_CAUSAL\_DECISION}.
\]
The replay authenticates finite algebra and strict finite signs.  The
asymptotic theorems are proved in the text, not inferred from the fixtures.

\section{Final boundary}

\begin{center}
\begin{tabular}{ll}
\toprule
Native first-owner coefficient ledger & proved exact\\
Residual $r-2r^2$ and $(R-A)F$ & proved exact\\
Root equals native annular scalar & proved exact\\
Raw-exposure conservation & proved exact\\
Subcritical threshold-depth positivity & false\\
$\LAPBR$ & false\\
High-least-prime sector & proved positive\\
$\CSHT$ & open / RH-bearing\\
Riemann Hypothesis & unproved\\
\bottomrule
\end{tabular}
\end{center}

\end{document}
